\section{Results and Discussion}
\label{sec:results}

\subsection{Landscape Graph Construction}
\label{sec:som-results}

The SOM converged after approximately 150 training epochs on the
400\,000-pixel subsample (80\,\% of 500\,000, with 20\,\% held out
for validation). \Cref{fig:som-convergence} shows the evolution of
validation quantisation error (QE) and topographic error (TE) across
training stages. QE decreases monotonically throughout training,
but the codebook feature maps
(\cref{fig:codebook-comparison}) reveal that this continued
improvement comes at a cost: at 150~epochs, the codebook displays
smooth gradients with distinct regional structure, whereas at
200~epochs the features fragment into noise as the shrinking
neighbourhood radius causes individual neurons to overfit local
samples. While QE continues to decrease beyond 150~epochs, the
topographic error increases and the codebook loses the spatial
coherence needed for meaningful landscape-type clustering. The
final validation metrics (QE\,=\,1.54, TE\,=\,0.14) at 150~epochs
thus represent the best trade-off between quantisation fidelity and
topological preservation.

\begin{figure}[htbp]
  \centering
  \includegraphics[width=0.75\linewidth]{som_convergence}
  \caption{SOM convergence on the held-out validation set. Left
  axis: quantisation error (mean $L_2$ distance to best-matching
  unit). Right axis: topographic error (fraction of samples whose
  two closest BMUs are non-adjacent on the grid). The dashed
  vertical line marks the 150-epoch production threshold.}
  \label{fig:som-convergence}
\end{figure}

Hierarchical clustering of the 4\,096 codebook vectors (Ward's
linkage, cut at 30 clusters) partitioned the SOM grid into 30
landscape types representing recurring environmental configurations
across Switzerland. Connected-component labelling and merging of
fragments smaller than 100~pixels produced the final landscape graph
of \textbf{166\,464 patches} and \textbf{465\,833 edges}.

The patch size distribution (\cref{fig:patch-stats}) shows a median
of 340~pixels (${\sim}0.31$\,km$^2$), with a long tail towards
larger contiguous patches in homogeneous terrain (e.g.\ alpine
grassland, dense forest). The minimum enforced size of 100~pixels
(${\sim}0.09$\,km$^2$) ensures that each patch is large enough to
carry ecologically meaningful feature averages.

The full-extent PCA-coloured patch map (\cref{fig:pca-patches}) reveals
coherent biogeographic structure: similar colours emerge for the
Swiss Plateau and the Po Valley south of the Alps, reflecting
analogous lowland environmental profiles. High-alpine zones form a
distinct cluster, while the southern Pre-Alps and Ticino group with
the Jura arc, both of which are characterised by their southward slope and
intermediate elevations.

\Cref{fig:landscape-roi} zooms into a ${\sim}16 \times 16$\,km
region north of Bern. Urban patches (Bern city centre) are clearly
delineated. The Dählhölzli forest along the Aare, the Schärmenwald
near Ittigen, and the Tiefenau peninsula emerge as distinct
forest-type patches. Overall, when compared with local knowledge,
clear ecological structures are visible at this resolution.

\begin{figure}[htbp]
  \centering
  \includegraphics[width=0.85\linewidth]{patches_bern}
  \caption{Landscape-type map for the Bern region. Colours
  represent the 30 landscape types derived from hierarchical
  clustering of the SOM codebook. The tessellation captures
  transitions between forested hillsides, agricultural plains,
  and built-up areas at 30\,m resolution.}
  \label{fig:landscape-roi}
\end{figure}

\paragraph{Discussion.}
The SOM-based discretisation serves a dual purpose: it reduces the
${\sim}$160~million valid pixels to a tractable graph of 166\,464
nodes while ensuring that patch boundaries follow ecological
gradients rather than an arbitrary regular grid. This adaptive
resolution means that patches are smaller in heterogeneous landscapes
(e.g.\ valley floors with interspersed land cover) and larger in
homogeneous terrain (e.g.\ high-alpine meadows), mirroring the
spatial structure relevant to species distributions.

A limitation of the adjacency-based edge construction is that
connections between patches sharing only a single boundary pixel
receive the same weight as connections along extensive shared
borders. The graph therefore contains ``bridge'' edges representing
tenuous spatial contact (e.g.\ two forest patches touching through a
one-pixel corridor under a road). This is revisited in
\cref{sec:limit-detail}.

\subsection{Random Forest Baseline}
\label{sec:rf-results}

A per-species Random Forest served as the non-spatial baseline.
Across 3\,751 species, the RF achieved a mean AUC of 0.838
(median\,=\,0.834), with 99.4\,\% of species exceeding 0.7. Full
metric summaries are provided in \cref{app:rf-baseline}. The RF's
key limitation is its independence assumption: each patch is
classified without regard to its landscape context, so the model
cannot distinguish isolated habitat fragments from patches embedded
in connected corridors.

\subsection{Architecture Search}
\label{sec:arch-results}

Six GraphSAGE architectures were evaluated on six
representative species spanning a range of habitats and
prevalence levels (\cref{tab:arch-search}).  Performance was
measured by validation AUC averaged over the six species.

\begin{table}[htbp]
  \centering
  \caption{Architecture search results (validation AUC). Each
  column corresponds to one of the six representative species;
  the final column reports the mean across species.}
  \label{tab:arch-search}
  \small
  \begin{tabular}{lcccccc|c}
    \toprule
    Architecture & \textit{P.~ab.} & \textit{F.~sy.}
      & \textit{T.~pr.} & \textit{P.~er.}
      & \textit{Ph.~au.} & \textit{S.~in.} & Mean \\
    \midrule
    1-layer {[}16{]}        & 0.730 & 0.755 & 0.670
      & 0.700 & 0.734 & 0.877 & 0.744 \\
    2-layer {[}24,12{]}     & 0.742 & 0.785 & 0.715
      & 0.725 & 0.746 & 0.885 & 0.766 \\
    paper {[}24,18,8{]}     & 0.770 & 0.791 & 0.726
      & 0.751 & 0.771 & 0.901 & 0.785 \\
    wider {[}48,32,16{]}    & 0.791 & 0.807 & 0.752
      & 0.766 & 0.791 & 0.908 & 0.803 \\
    extra-wide {[}64,48,32{]} & 0.796 & 0.817 & 0.759
      & 0.776 & 0.794 & 0.911 & 0.809 \\
    4-layer {[}32,24,16,8{]} & 0.787 & 0.796 & 0.742
      & 0.750 & 0.776 & 0.906 & 0.793 \\
    \bottomrule
  \end{tabular}
\end{table}

The extra-wide architecture ({[}64,48,32{]}, 3~layers) yielded the
highest mean AUC (0.809) and was selected for all subsequent
experiments.  Adding a fourth layer did not improve performance,
suggesting that three message-passing hops (covering a
${\sim}$3-patch spatial neighbourhood) capture the relevant
landscape context for Swiss plant species.  This contrasts with
the original \cite{Wu2025} configuration ({[}24,18,8{]}),
which was optimised for global-scale fauna; wider hidden layers
appear necessary to capture the finer environmental gradients
that distinguish Swiss vegetation zones at 30\,m resolution.

\paragraph{Graph structure benefit.}
The improvement from one to three GraphSAGE layers quantifies how
much each species profits from neighbourhood information
propagated through the landscape graph. The largest gains are
observed for alpine and subalpine grassland species---notably
\textit{Potentilla erecta} ($+0.051$ AUC) and \textit{Trifolium
pratense} ($+0.056$)---while the habitat specialist
\textit{Senecio inaequidens} benefits least ($+0.024$), as its
distinct urban niche is already well captured by local features
alone.

\textit{Potentilla erecta} (alpine grassland) provides a
particularly clear illustration. This species occupies connected
subalpine meadow patches in distinct elevation bands, shaped by
climate and topography rather than human management. With a
single GraphSAGE layer, predictions rely solely on local patch
features; isolated meadow patches receive unreliable scores
because the model lacks spatial context. With three layers, the
GNN propagates suitability information through adjacent patches,
correctly identifying connected alpine grassland belts as
suitable habitat---even where direct observations are absent.
This spatial smoothing through the graph is the core mechanism
by which SOM-GNN-SDM surpasses the non-spatial RF baseline for
spatially structured species (see \cref{sec:gnn-results}).

\subsection{SOM-GNN-SDM Performance}
\label{sec:gnn-results}

% Same metrics, GNN vs RF comparison
% Where does graph structure help?
% Where does RF suffice (generalists)?

\textbf{TODO: GNN results and comparison with RF.}


\subsection{Common vs.\ Vulnerable Species}
\label{sec:species-comparison}

% Deep dive on 12 species (6 common + 6 vulnerable)
% Overprediction discussion: sampling bias, SDM purpose
% Specialists vs generalists performance difference

\textbf{TODO: Species-level analysis.}

\subsection{Habitat Suitability Maps}
\label{sec:maps}

% Example maps for selected species
% Visual validation against known ranges

\textbf{TODO: Map figures.}

\subsection{Conservation Prioritisation}
\label{sec:conservation-results}

% Circuit-theory connectivity results
% Gap analysis: priority areas vs protected areas
% Quantification of conservation gaps

\textbf{TODO: Conservation application results.}
